% !TEX root = ../main.tex
% =====================================================================
% 5. RESULTS -- page budget 2.60 (three floats live here)
%
% One subsection per research question, in the order the reader asks
% them: does it still work, does it still explain, what did it cost,
% which component did what, and can the faithfulness numbers be trusted.
% =====================================================================

\section{Results}
\label{sec:results}

\begin{table}[t]
\caption{Main comparison. \webtraffic{} has per-time-step ground truth; \ucr{}-85 is the wider check
and was never used to choose a model. $S$ = seeds; rows with more than one seed are means. Cost is at
batch 32, length 1008, 10 classes. $^{\ddagger}$84 of 85 datasets finished.
$^{\dagger}$\millet{}'s own 1{,}500-epoch recipe: a harder accuracy target, and an AOPCR value that
\emph{cannot} be compared with the rest of that column (\Cref{sec:res:aopcr}). Bold marks the best
matched-budget row in each column.}
\label{tab:main}
\begin{center}
\footnotesize
\setlength{\tabcolsep}{3.2pt}
\begin{tabular}{llccccrr}
\toprule
& & \multicolumn{3}{c}{\webtraffic{}} & \ucr{}-85 & \multicolumn{2}{c}{Cost} \\
\cmidrule(lr){3-5}\cmidrule(lr){6-6}\cmidrule(lr){7-8}
\textbf{Model} & $S$ & Acc. $\uparrow$ & AOPCR $\uparrow$ & \ndcg{} $\uparrow$ & Acc. $\uparrow$ &
Params & MFLOPs \\
\midrule
FCN $+$ conjunctive           & 1 & 0.742 & 3.826 & 0.533 & 0.814 & 267{,}035 & 535.2 \\
ResNet $+$ conjunctive        & 1 & 0.772 & 2.952 & 0.554 & 0.815 & 506{,}331 & 1013.2 \\
\millet{} (matched budget)    & 3 & 0.887 & 2.569 & 0.677 & \textbf{0.827}$^{\ddagger}$ & 423{,}707 & 847.6 \\
\millet{} (published recipe)$^{\dagger}$ & 1 & 0.920 & 13.268 & 0.661 & 0.843$^{\ddagger}$ & 423{,}707 & 847.6 \\
\midrule
\seanet{} wide $+$ class-wise & 1 & 0.946 & 1.722 & 0.686 & \textbf{0.829} & 269{,}164 & 523.7 \\
\seanet{} gated $+$ Top-$k$   & 3 & \textbf{0.955} & 2.225 & 0.750 & 0.808 & 61{,}740 & 109.7 \\
\textbf{\seanet{} bottleneck $+$ Top-$k$ (ours)}
                              & 3 & 0.947 & \textbf{2.621} & \textbf{0.772} & 0.810 & \textbf{41{,}324} & \textbf{76.7} \\
\bottomrule
\end{tabular}
\end{center}
\end{table}

\subsection{Classification and Interpretation on WebTraffic}
\label{sec:res:main}

\Cref{tab:main} answers RQ1 and RQ2 on \webtraffic{}, and both point the same way. Against the
baseline trained under our budget, the 41\,K model is $6.0$ accuracy points better ($0.947$ against
$0.887$, three seeds each, seed standard deviation $\pm 0.016$) and $9.5$ \ndcg{} points better
($0.772$ against $0.677$, $\pm 0.013$); both gaps are several times the seed spread. AOPCR is $2.621$
against $2.569$, which is inside the AOPCR standard deviation of either model ($\pm 0.214$ and
$\pm 0.884$), so it should be read as a tie. In short: better classification, better interpretation
ranking, and no change in faithfulness --- with $10.3\times$ fewer parameters.

Two rows put that in context. \millet{} under its own 1{,}500-epoch recipe reaches $0.920$, so more
than half of our gain over the matched-budget baseline comes from the training budget and not from
the architecture; our model is still ahead of it by $2.7$ accuracy and $11.1$ \ndcg{} points at
$10.3\times$ less cost. And \seanet{} gated $+$ Top-$k$, at $61{,}740$ parameters, is on paper the
most accurate row at $0.955$ --- but $0.008$ against a seed spread of $0.016$ is not a difference we
claim, and we prefer the bottleneck version because it is $33\,\%$ smaller with no loss on \ndcg{}.

\subsection{Model Cost}
\label{sec:res:cost}

\begin{table}[t]
\caption{Cost, measured at batch 32, length 1008. Weight memory is \emph{calculated} from the
parameter count (4\,B and 1\,B per weight); nothing was quantised or run on a device. Latency and
peak memory are GPU measurements, not estimates for a device.}
\label{tab:cost}
\begin{center}
\small
\setlength{\tabcolsep}{5pt}
\begin{tabular}{lrrrrrr}
\toprule
\textbf{Model} & Params & \multicolumn{2}{c}{Weights (KB)} & MFLOPs & Lat. & Mem. \\
\cmidrule(lr){3-4}
 & & fp32 & int8 & & (ms) & (MB) \\
\midrule
\millet{}               & 423{,}707 & 1655 & 414 & 847.6 & 0.203 & 112.3 \\
\seanet{} wide          & 269{,}164 & 1051 & 263 & 523.7 & 0.359 & 123.7 \\
\seanet{} gated         &  61{,}740 &  241 &  60 & 109.7 & 0.128 &  64.9 \\
\textbf{\seanet{} bottl.} & \textbf{41{,}324} & \textbf{161} & \textbf{40} & \textbf{76.7} & 0.131 & 179.0 \\
\seanet{} bottl., 2 blocks & 21{,}548 &   84 &  21 &  40.0 & 0.069 & 178.9 \\
\bottomrule
\end{tabular}
\end{center}
\end{table}

\Cref{tab:cost} separates what improved from what did not. Parameters fall $10.3\times$ and
computation $11.1\times$; at 32-bit precision the weights go from about $1.6$\,MB to $161$\,KB, which
moves the model from above to well below the roughly 1\,MB flash budget of the microcontroller class
described by \citet{lin2020mcunet} and \citet{banbury2021micronets}, and standard int8 quantisation
\citep{jacob2018quantization} would take it to about $40$\,KB. That is the sense in which it becomes
a candidate for such a device.

The other two columns do not improve, and we report them instead of leaving them out. Latency
improves by only $1.6\times$ even though the FLOPs fall $11.1\times$, because depthwise convolutions
are limited by memory speed and a GPU at batch 32 is the wrong tool for a claim about small hardware
--- the $269$\,K wide model is in fact the \emph{slowest} row. Peak memory use is $179.0$\,MB against
the baseline's $112.3$\,MB, which is worse, and it is worse because of the design: the encoder keeps
the length unchanged, so the activations are $(\dmodel, \nsteps)$ at every stage no matter how few
weights produce them, and \Cref{eq:bottleneck} reduces weights, not activations. On microcontrollers
the SRAM needed for activations is the limit at least as often as the flash needed for weights, so
this is a real limitation of the design and not a measurement error (\Cref{sec:disc}).

\subsection{Accuracy on the UCR Archive}
\label{sec:res:ucr}

\ucr{} was never used to choose anything, and it is where the trade-off shows. Our 41\,K model
reaches $0.810$ mean accuracy against $0.827$ for the matched-budget baseline and $0.843$ for
\millet{}'s published recipe: $1.8$ and $3.4$ points behind, with the same order in mean rank over
the 84 shared datasets, $15.40$ against $12.60$ and $9.41$ (\Cref{app:ucr}).

Where that loss comes from matters more than its size, and one row of \Cref{tab:main} answers it. The
\emph{wide} version of the same encoder --- same blocks and same kind of head, with $\dmodel$ raised
from 64 to 128 and 4 blocks to 6 --- reaches $0.829$ and mean rank $11.53$, ahead of the
matched-budget baseline on both. So the architecture is not what loses on \ucr{}; the smaller size is.
Two comparisons follow, and they are the useful numbers here: changing the backbone at a similar size
is worth about $1.1$ rank places, while going from our 400-epoch budget to \millet{}'s 1{,}500-epoch
budget on a \emph{fixed} architecture is worth about $3.2$ --- three times more.

\subsection{Ablation Study}
\label{sec:res:ablation}

\begin{table}[t]
\caption{Ablations on \webtraffic{}, all at $\dmodel = 64$, 4 blocks, seed 0. \emph{Left:} the two
parts on and off --- reading down gives the bottleneck, reading across gives \mbox{Top-$k$}.
\emph{Right:} only $\kappa$ changes, and at $\kappa = 1$ the head \emph{is} the class-wise baseline
(\Cref{prop:topk}); the last row instead turns off the entropy term of \Cref{eq:loss}.}
\label{tab:ablation}
\begin{center}
\begin{minipage}[t]{0.51\textwidth}
\centering
\scriptsize
\setlength{\tabcolsep}{3.5pt}
\begin{tabular}{llrccc}
\toprule
\textbf{Encoder} & \textbf{Head} & Params & Acc. & AOPCR & \ndcg{} \\
\midrule
plain      & class-wise & 57{,}580 & 0.916 & 2.478 & 0.719 \\
plain      & Top-$k$    & 57{,}580 & 0.932 & 2.691 & \textbf{0.778} \\
bottleneck & class-wise & \textbf{41{,}324} & 0.902 & 2.572 & 0.733 \\
bottleneck & Top-$k$    & \textbf{41{,}324} & \textbf{0.938} & \textbf{2.778} & 0.777 \\
\bottomrule
\end{tabular}
\end{minipage}\hfill
\begin{minipage}[t]{0.45\textwidth}
\centering
\scriptsize
\setlength{\tabcolsep}{4pt}
\begin{tabular}{lccc}
\toprule
$\kappa$ & Acc. & AOPCR & \ndcg{} \\
\midrule
0.05 & 0.888 & 2.288 & 0.715 \\
\textbf{0.10} & 0.938 & 2.778 & \textbf{0.777} \\
0.25 & \textbf{0.940} & 2.648 & 0.733 \\
0.50 & 0.882 & \textbf{3.114} & 0.732 \\
1.00 \emph{(= baseline head)} & 0.908 & 2.436 & 0.722 \\
\midrule
$\lambda_{\mathrm{focus}} = 0$ & 0.950 & 2.303 & 0.765 \\
\bottomrule
\end{tabular}
\end{minipage}
\end{center}
\end{table}

\textbf{RQ4, the bottleneck.} Reading down \Cref{tab:ablation}, it removes $16{,}256$ parameters ---
$28\,\%$ of the model --- and accuracy moves by $-0.014$ with the class-wise head and $+0.006$ with
\mbox{Top-$k$}, both inside the $\pm 0.016$ seed spread; \ndcg{} moves by $+0.014$ and $-0.001$. So
the safe statement is that the bottleneck buys the size reduction \emph{with no accuracy cost we can
measure}, not that it improves anything: the rest of the model does not notice the change.

\textbf{RQ5, \mbox{Top-$k$} pooling.} Reading across, \ndcg{} rises by $0.059$ on the plain encoder
and $0.044$ on the bottleneck encoder. Against a \ndcg{} seed spread of $\pm 0.013$ these are three
to four standard deviations, and they are the claim we make for the head; the accuracy changes
($+0.016$ and $+0.036$) point the same way but are at or below one seed spread, so we do not present
\mbox{Top-$k$} as an accuracy improvement. The right panel of \Cref{tab:ablation} tests the selection
on its own, and here \Cref{prop:topk} earns its place: the $\kappa = 1$ row is exactly the baseline
head, so the $0.777$ against $0.722$ gap between $\kappa = 0.1$ and $\kappa = 1$ is caused by the
selection alone. Two things in that table are worth saying plainly instead of hiding. There is a
\emph{best value}, not a steady gain --- $\kappa = 0.05$ is worse than $\kappa = 0.10$ in every
column, so keeping too few steps is as harmful as keeping all of them, and $\kappa$ has to be chosen.
And AOPCR does not follow $\kappa$ at all: it is highest at $\kappa = 0.50$, where \ndcg{} is near its
lowest. The last row shows that the attention-entropy term is a second, separate control that works in
the same direction: turning it off improves accuracy ($0.950$) but costs \ndcg{} ($0.765$) and AOPCR
($2.303$). Every row is a single seed, so we read the table for the overall pattern and not for single
values.

\subsection{Sensitivity of AOPCR to the Training Budget}
\label{sec:res:aopcr}

One number in \Cref{tab:main} needs an explanation, and the explanation applies more widely.
\millet{} under its published recipe scores AOPCR $13.268$ on \webtraffic{}, five times every other
row, with the same architecture, code, metric implementation and data as the matched-budget
\millet{} row, which scores $2.569$; only the training recipe is different. On \ucr{}-85 the same
pair gives $0.812$ against $3.926$, a factor of $4.84$. The reason is that AOPCR adds up a drop in
the logit that has \emph{no fixed scale} (\Cref{app:metrics}): a model trained longer and to a lower
loss has sharper logits, so changing the input produces a larger absolute drop whether or not the
explanation points at the right place. In line with this, the same pair moves by only $1.6$ points of
\ndcg{} --- the measure that does have a fixed scale --- and moves \emph{down}, $0.677$ to $0.661$,
while AOPCR is multiplied by $5.2$. So an AOPCR value from one paper cannot be compared with one from
another: it should only be reported between models trained the same way, or replaced by a measure
with a fixed scale wherever ground truth allows. That is why we retrain every baseline instead of
copying published numbers, and never compare against \millet{}'s published AOPCR --- including where
it would make our results look better.
